\begin{theorem}[Trace pairings across a cut]%
    \label{thm:w_state_one_sided_cut_rank}
    \lean{MPSTensor.blockTracePairing_range_le_of_forall_mem,
        MPSTensor.sq_dim_le_finrank_of_forall_mem}
    \leanok
    \uses{def:l_blk_injective}
    Let $B^1,\ldots,B^b$ have bond dimensions $D_1,\ldots,D_b$ and suppose
    their word tuples of length $R'$ span $\bigoplus_j\MN{D_j}$. Let every
    $c_j\ne0$. If each function
    $\sigma\mapsto\sum_jc_j\tr(B^j_\sigma B^j_\tau)$ lies in a subspace $V$
    of functions on words of length $R$, then every function
    $\sigma\mapsto\sum_j\tr(X_jB^j_\sigma)$ also lies in $V$.
    If the specified block $B^j$ is $R$-block injective, then $D_j^2\le\dim V$.
\end{theorem}

\begin{proof}
    \leanok
    For an arbitrary tuple $(X_j)_j$, simultaneous spanning gives coefficients
    $a_\tau$ such that $\sum_\tau a_\tau c_jB^j_\tau=X_j$. Hence
    \begin{align}
        \sum_\tau a_\tau\sum_jc_j\tr(B^j_\sigma B^j_\tau)
         & =\sum_j\tr\left(B^j_\sigma\sum_\tau a_\tau c_jB^j_\tau\right)
        =\sum_j\tr(X_jB^j_\sigma).
        \label{eq:w_state_trace_range}
    \end{align}
    To obtain the dimension bound, choose $X_k=0$ for $k\ne j$.
    The map $X_j\mapsto(\tr(X_jB^j_\sigma))_\sigma$ is injective: if all its
    values vanish, $R$-block injectivity gives $\tr(X_jY)=0$ for every matrix
    $Y$, and matrix units then give $X_j=0$. Its image lies in $V$ by
    (\ref{eq:w_state_trace_range}), so $D_j^2\le\dim V$.
\end{proof}

\begin{theorem}[W-state bound from unequal cuts]%
    \label{thm:w_state_asymmetric_bound}
    \lean{MPSTensor.lt_of_sum_mpv_eq_smul_wIndicator_asymmetric,
        MPSTensor.PGVWC07CanonicalFormData.lt_of_mpv_eq_smul_wIndicator_asymmetric}
    \leanok
    \uses{def:pgvwc07_ti_canonical_representation, def:l_blk_injective,
        def:blocks_not_gauge_phase_equiv, def:w_indicator}
    Let $B^1,\ldots,B^b$ be binary tensors of positive bond dimensions, each
    unital and irreducible, with a positive definite fixed point of its dual
    transfer map. Suppose no two blocks differ by a gauge transformation and
    a phase, and each is $L_0$-block injective for a common $L_0>0$.
    If $\sum_jc_j\ket{V^{(N)}(B^j)}=c\ket{W_N}$ with every $c_j\ne0$ and
    $c\ne0$, then, writing $F=\max\{L_0,3(b-1)(L_0+1)\}$,
    \begin{align}
        N & <\max\{L_0+F,12(b-1)\}.
        \label{eq:w_state_asymmetric_bound}
    \end{align}
    The same conclusion holds for canonical-form coefficients $c_j=\lambda_j^N$.
    Here $b-1$ means zero when $b=0$.
    The threshold in~(\ref{eq:w_state_asymmetric_bound}) is at most $2F$;
    for $L_0\ge3$ it is $L_0+F$.
\end{theorem}

\begin{proof}
    \leanok
    \uses{thm:w_state_one_sided_cut_rank, lem:w_state_cut_rank,
        lem:pgvwc07_block_tuple_span, thm:w_state_canonical_bound}
    Suppose $N\ge\max\{L_0+F,12(b-1)\}$ and cut at $L_0$.
    The word tuples on the other side, of length $N-L_0\ge F$, span the
    product algebra by Lemma~\ref{lem:pgvwc07_block_tuple_span}.
    Individual injectivity on the first side and
    Theorem~\ref{thm:w_state_one_sided_cut_rank} give $D_j^2\le2$, since
    the W-state cut rank is at most two. Thus every $D_j=1$.

    For scalar letters $b_i^j$, unitality reads $\sum_i|b_i^j|^2=1$.
    A nonzero letter spans the one-dimensional matrix algebra, so the
    same blocks are injective at length one. Reapply
    Theorem~\ref{thm:w_state_canonical_bound} with this length to obtain
    \begin{align}
        N & <2\max\{1,6(b-1)\}
        =\max\{2,12(b-1)\}
        \le\max\{L_0+F,12(b-1)\},
    \end{align}
    a contradiction. Finally, $L_0+F\le2F$ and $12(b-1)\le2F$.
    If $L_0\ge3$ and $b\ge2$, then $3(b-1)(L_0+1)\ge12(b-1)$;
    if $b\le1$, the latter term is zero.
\end{proof}

\begin{theorem}[Cubic bound from a nonzero one-step eigenvalue]%
    \label{thm:w_state_nonzero_eigenvalue_bound}
    \lean{MPSTensor.PGVWC07CanonicalFormData.lt_of_mpv_eq_smul_wIndicator_of_nonzero_eigenvalue}
    \leanok
    \uses{def:pgvwc07_ti_canonical_representation, def:l_blk_injective,
        def:blocks_not_gauge_phase_equiv, def:w_indicator}
    Let $A$ be a binary tensor of bond dimension $D$ in canonical form, with
    no two blocks related by a gauge transformation and a phase. Suppose each
    block $B^j$ is injective at some positive length and there are
    $X_j\in\spn_{\C}\{B_i^j\}$, $\mu_j\ne0$, and $\phi_j\ne0$ satisfying
    $X_j\phi_j=\mu_j\phi_j$. If $V^{(N)}(A)=cW_N$ with $c\ne0$, then,
    for $L=\max\{1,D^2\}$,
    \begin{align}
        N & <\max\bigl\{L+\max\{L,3(D-1)(L+1)\},12(D-1)\bigr\}.
        \label{eq:w_state_cubic_eigenvalue}
    \end{align}
    For positive $D$, $L=D^2$, so this gives $D=\Omega(N^{1/3})$.
    The maximum with one includes $D=N=0$.
\end{theorem}

\begin{proof}
    \leanok
    \uses{thm:w_state_asymmetric_bound, thm:wielandt_general}
    Positive-length injectivity implies full word span at every larger length.
    The nonzero-eigenvalue case of the quantum Wielandt bound
    \cite[Theorem~6.9]{Wolf2012Quantum} gives injectivity at $D_j^2$ in each block.
    Since $D_j\le D$, this propagates to the common length $L$.
    Apply Theorem~\ref{thm:w_state_asymmetric_bound} and use $b\le D$.
\end{proof}
